\scanpage{66}
\mspage{65}
\mspage{33}

La suite spectrale $(R^pv^!)(R^qu^!)F\Longrightarrow(R^{p+q}w^!)F$ donne un isomorphisme
\diagram{\[
\xi:\begin{tikzcd}
(R^{2b}v^!)(R^{2a}u^!)\mu_n^{\otimes c}\arrow[r,"\sim"]&(R^{2c}w^!)\mu_n^{\otimes c}
\end{tikzcd}.
\]}
On obtient ainsi un \unsure{diagbamme|diagramme} d'isomorphismes
\diagram{\[
(3.10.1)\quad
\begin{tikzcd}[column sep=large,row sep=large]
\mathbb Z/n\arrow[r,"{\varphi(Z,Y)}"]\arrow[d,"{\varphi(Z,X)}"]&(R^{2b}v^!)\mu_n^{\otimes b}\arrow[d,"\eta="',"{(R^{2b}v^!)(\ (Y,X)\otimes\mu_n^{\otimes b})}"]\\
(R^{2c}w^!)\mu_n^{\otimes c}&(R^{2b}v^!)(R^{2a}u^!)\mu_n^{\otimes c}\arrow[l,"\xi"]
\end{tikzcd},
\]}
et l'assertion de transitivité est que cet diagramme est commutatif.

Démonstration. La démonstration de l'existence et unicité se faite par recurrence sur la codimension. Nous commençons avec $c=1$, le cas $c=0$ étant trivial. Notons que le morphisme $\varepsilon_t$ de (3.9) est de nature locale sur $X$ pour la topologie étale et commute bien aux changements de base, avec la convention évident \unsure{pout|pour} $t$. Puisque le transitivité (iii) n'intervient guère pour $c=1$, tout revient à démontrer que (3.9) est bijectif et indépendant \unsure{au|du} choix de $t$. Or cela\deleted{\unsure{n}} se vérifie fibre par fibre sur $Y$, et puisque (3.9) et $(R^qi^!)$ (3.8) commutent aux changements de base, on peut supposer $S=\operatorname{Spec}k$, où $k$ est un corps séparablement clos
